\documentclass[a4paper]{article}
\usepackage[pdftex]{graphicx}
\usepackage[utf8]{inputenc}
\usepackage{enumerate}
\usepackage{icomma}
\usepackage{amssymb}
\usepackage{tikz}
\usepackage{href-ul}
\hypersetup{
	colorlinks=true,
	linkcolor=blue,
	urlcolor=blue}
\usepackage{geometry}
\geometry{a4paper, top=15mm, left=15mm, right=15mm, bottom=15mm,
	headsep=10mm, footskip=12mm}

\begin{document}
	\begin{tikzpicture}(10,3)
		\draw[ultra thick](2,0) --(10,0) -- (10,1.3) --(2,1.3) -- (2,0);
		\draw[fill=black](2,0)-- (0,.6) -- (2,.6) -- (2,0);
		\node at (6,.6) {\large solution sheet from \href{http://www.okuyakl.com}{www.okuyakl.com}};
	\end{tikzpicture}
	\vspace{0.5 cm}
	
	\noindent {\bf task 1. a)} \\
	First derivative with chain rule:
	$$f'(x)=-2\ln{x}\cdot{1 \over x}$$
	Second derivative with product rule:
	$$
	\begin{array}{rcl}
		f''(x)&=&{2 \over x^2} \cdot \ln{x} - {2 \over x} \cdot {1\over x}\\
		\ &=& {2 \over x^2} \cdot \left(\ln{x}-1 \right)\stackrel{!}{=}0\\
		\ln{x_0}&=&1 \\
		x_0&=&e
	\end{array}
	$$
	
	Tangent: $$y_T=m \cdot x+t; \qquad x=e; \qquad m=f'(e)= -{2 \over e}; \qquad y_T= f(e)=0 $$
	$$ 0=-{2 \over e} \cdot e + t; \quad \Rightarrow t=2$$
	$$ \Rightarrow y= -{2 \over e}\cdot x +2 $$
	
	\noindent {\bf task 1. b)} \\
	Derive $F$ with the product and chain rule:
	$$
	\begin{array}{rcl}
		F'(x)&=&-(\ln{x}-1)^2 - x \cdot 2 \cdot (\ln{x}-1) \cdot {1 \over x}\\
		\ &=& -((\ln{x})^2 - 2 \ln{x}+1) - 2 \ln{x} + 2 \\
		\ &=& 1-(\ln{x})^2
	\end{array}
	$$
	
	\noindent {\bf task 1. c)}\\
	Zeropoint:
	$$
	\begin{array}{rcl}
		1-(\ln{x})^2 &=& 0 \\
		\ln{x} &=& \pm 1 \\
		x_1&=&e \\
		x_2&=&e^{-1}
	\end{array}
	$$
	With this and the given antiderivative we calculate the area:
	$$A=\int\limits_{e^{-1}}^e f(x) ~{\rm d} x =\left[ -x(\ln x -1)^2\right]_{e^ {-1}}^e = 0 -(-e^{-1}(-2)^2) = {4\over e}$$
	
	\noindent {\bf task 1. d)}\\
	Every integral function has a zero point that lies where the upper limit, here the variable $x$, coincides with the lower limit, here the 0, i.e. at $x=0$. Then $I(0)=0$, here is
	also $F(0)=0$, so $F(x)=I(x)$
	
	\noindent {\bf task 1. e)}\\
	According to the mnemonic, the limit is $$\lim \limits_{x \to 0^+} x \cdot ln{x} =0 $$
	and therefore also:
	$$\lim \limits_{x \to 0^+} -x \cdot ( ln{x}-1)^2 =0 $$
	This means that the area bounded by the function and the two axes in the fourth quadrant is finite.
	
	\noindent {\bf task 2. a)}\\
	We have: $I'(x)=f(x)$, so we set $f(x)=0$:
	$$
	\begin{array}{rcl}
		(e^x-2)^2 &=& 0\\
		(e^x-2)&=& 0 \\
		e^x &=& 2 \\
		x&=& \ln{2}
	\end{array}
	$$
	Monotonicity table:
	$$
	\begin{array}{cc|c|c}
		& x< \ln{2} & x=\ln{2} & x> \ln{2} \\
		\hline
		(e^x-2)^2 & + & 0 & + \\
		I'(x) & + & 0 & + \\
		\hline
		I(x) & \nearrow & \rightarrow & \nearrow \\
		& sms & terrace point & sms
	\end{array}
	$$
	The function $I(x)$ is strictly increasing monotonically throughout $\mathbb{R}$ and has a terrace point at $T(\ln2 | 0 )$.
	
	\noindent {\bf task 2. b)}
	
	$$ F'(x)= 0.5 \cdot 2 \cdot e^{2x} -4e^x+4 = (e^x-2)^2$$
	
	\noindent {\bf task 2. c)}
	Determination of the upper integration limit:
	$$
	\begin{array}{rcl}
		(e^x-2)^2 &=& 4 \\
		e^x-2 &=&\pm 2 \quad \textnormal{ only positive solution possible}\\
		e^x&=&4 \\
		x &=& \ln{4}
	\end{array}
	$$
	$$
	\begin{array}{rcl}
		A(u)&=&\int \limits_u^{\ln{4}}(4- (e^x-2)^2 )~{\rm d} x = \left[ 4x-0.5\cdot e^{2x}+4e^x-4x \right]_u^{\ln{4}}\\
		A(u)&=&-0.5\cdot 16 + 4 \cdot 4 -0.5 \cdot e^{2u} - 4e^{u}\\
		&=& 8-0.5 \cdot e^{2u} - 4e^{u}
	\end{array}
	$$
	
	
	\noindent {\bf task 2. d)}
	$$\lim \limits_{u \to -\infty} 8-0.5 \cdot e^{2u} - 4e^{u}=8-0-0=8$$
	The area enclosed by the graph and the line $y=4$ has a finite area and is $8~FE$.
	
	\noindent {\bf task 3.}\\
	$I(x)$ is represented by graph {\bf B}. Graph {\bf A} has two terrace points, these would correspond to two double zeros of $G_f$. However, these are not available.
	Graph {\bf C} has 5 extrema, so it is at least 6th order, but only a function of order 3+1, i.e. 4, is possible.
	
	\begin{center}
		\includegraphics[width=7 cm]{../../viecher/eendcomic.pdf}
		
		Here you can go back to the \href{https://www.okuyakl.de/math/m12intapL222/ae222.pdf}{task sheet}
	\end{center}
	
\end{document}